\subsection{Analyse}

Dans son ensemble chaque experience en DSC est visible en Fig. \ref{fig:combined}.

\begin{figure}[hct]
        \centering
        \includegraphics[scale=0.45]{combined.pdf}
        \label{fig:combined}
        \caption{Vue d'ensemble des mesures DSC}
\end{figure}

\\
Avec la chauffe rapide en vert sapin, le refroidissement rapide en rouge, la chauffe lente en jaune et le refoidissement lent en violet (voir Fig. \ref{fig:combined})...
\\

Que l'on peut resumer dans le tableau chronologique comparatif suivant :
\begin{table}[!ht]
\centering
\footnotesize
\setlength{\tabcolsep}{4pt}
\renewcommand{\arraystretch}{0.9}
\begin{tabular}{|c|c|c|c|c|}
\hline
	\textbf{Condition (C|R)} &
\textbf{Cristallisation froide ?} &
\textbf{Fusion} &
\textbf{Cristallisation chaude ?} &
\textbf{Cristallinite} \\
\hline
C rapide       & ??  & ?? & ?? & ?? \\
R rapide & ??  & ?? & ?? & ?? \\
C lent         & ?? & ?? & ?? & ?? \\
R lent   & ?? & ?? & ?? & ?? \\
\hline
\end{tabular}
        \caption{tableau comparatif des mesures DSC}

\end{table}

